% Notes-only A.82 candidate. Unregistered; no canonical promotion.
\subsection{The full order-sixty-four wreath cover and its intermediates}

\begin{statement}{A.82}{Open}{The full wreath closure and every intermediate}
Let \(C\to E\) be the actual Galois closure of the cover in A.80,
with group \(G=(C_2)^4\rtimes C_4\), and let \(H\le G\).
Put \(Z=C/H\), let \(f_Z:Z\to E\) be the actual quotient map,
and put \(X_Z=Z\setminus f_Z^{-1}(O)\).
For every original Part A packet, every coefficient prime \(\ell\)
and every permitted open \(\Delta_0\) preserving this cover and
selected cusp \(P\), the proposed conclusion is
\[
 \bigl(\pi_1((X_Z)_{\overline L})^{(\ell)}\bigr)^{s_b(\Delta_0)}
                           =I_{P,\ell}.                 \tag{A.82.1}
\]
At odd \(\ell\), the same assertion holds for every cusp packet
over a finite number field and every open subgroup of that field's
closed derived Galois group. At two the assertion here is exactly
the original-packet scope of A.22. Every nontrivial closed procyclic
subgroup of the geometric pro-\(\ell\) group whose arithmetic
normalizer has open image in the absolute Galois group is contained
in a conjugate of cusp inertia, by A.59.
The assertion concerns this actual closure and every one of its
intermediates; no universal or cofinal family assertion is made.
\end{statement}

\begin{proof}[Complete notes-only candidate proof]
\paragraph{The actual character quotients.}
Retain all coordinates of A.80:
\[
 E_4:y^2=x^3-6x^2+x,\quad
 R=(1,2i),\quad T=2R=(0,0),\quad 2Q=R,
 \qquad K=\{O,R,T,-R\}.
\]
The degree-four isogeny \(\eta:E_4\to E\) has kernel \(K\).
Let \(a=(x-1)/\ell_Q\), with \(\ell_Q\) the tangent line
at \(Q\). Its divisor is \(R+O-2Q\).
The actual full function field is obtained by adjoining the square
roots of \(a_j(P)=a(P+jR)\), \(0\le j<4\).
A.80 proves their independence, the literal cyclic permutation
lift of translation by \(R\), and the resulting wreath group.
Write \(N=(C_2)^4\). For a nonempty subset \(S\) of the four
indices the quotient \(Y_S=C/\ker\chi_S\) is the actual double
cover of \(E_4\) defined by \(\prod_{j\in S}a_j\).
Its branch divisor is the symmetric difference of the indicated
pairs \(\{-jR,R-jR\}\).

For the rational \(N\)-character decomposition, the trivial
character contributes \(J_{E_4}\). Every nontrivial character
contributes the Prym of the indicated actual double cover.
Indeed the pullback and norm composition is
\(q_S^*q_{S,*}=\sum_{g\in\ker\chi_S}g\); subtracting the
trivial projector isolates the character. The orthogonal character
idempotents sum to one on rational Tate homology. The corresponding
sum of actual pullback maps is an isogeny onto \(J_C\), after
retaining each Prym rather than repeatedly retaining the base
elliptic factor. The four cyclic subset types give dimensions
\[
 1+4\cdot1+4\cdot1+4\cdot1+2\cdot2=17;              \tag{A.82.2}
\]
the all-four subset is an unramified elliptic cover and has Prym
zero. These are actual quotient projectors, not a decomposition
inferred from the genus alone.

The singleton Pryms are the elliptic \(F\) of A.80. The triple
Pryms are its conjugate \(F'\). Here is an actual square-class
check for the latter claim. Put \(D_K=\sum_{U\in K}U\) and
\(D_{Q+K}=\sum_{U\in K}(Q+U)\). Then
\[
 \operatorname{div}\prod_j a_j=2D_K-2D_{Q+K},
 \qquad [D_K-D_{Q+K}]=T\in\operatorname{Pic}^0(E_4).
                                                               \tag{A.82.3}
\]
The complementary triple to \(a_0\) differs by a square from
the double-cover function with divisor \(R+O-2(Q+T)\).
The half-difference of the divisors has degree zero and class
\(T-R+Q+(Q+T)=0\). On an elliptic curve such a divisor is
principal. The residual constant becomes a square after a finite
extension. Translation by \(-(Q+T)\) then gives the same
sextic model as A.80 with
\[
 q'=x(Q+T)=q^{-1}=3-2s+2s^2=q(-s).                    \tag{A.82.4}
\]
Thus this is the stated geometric conjugate, through actual
function-field isomorphisms. It is non-CM by A.80.

For an adjacent pair, say \(a_0a_1\), its divisor is
\(2O+R+(-R)-2Q-2(-Q)\). It equals the divisor of
\((x-1)/(x-q)^2\). Removing the square denominator gives the
actual genus-two model
\[
 v^2=(z^2+1)^3-6(z^2+1)^2+(z^2+1).
\]
The same two involution quotients as in A.80 give its Prym
\[
             F_{10976}:y^2=x^3-8x^2+12x+16,
                         \qquad j=10976.              \tag{A.82.5}
\]

\paragraph{The opposite pair, with the isogeny pullback identified.}
Let \(\phi:E_4\to E\) be the actual degree-two map of A.80,
\[
 \phi(x,y)=\left(\frac{(x-1)^2}{4x},
                        \frac{y(1-x^2)}{8x^2}\right).
\]
Put \(Q_0=\phi(Q)\), \(T_0=\phi(R)=(0,0)\). The exact
coordinates are
\[
           q_0=x(Q_0)=s^2+1=-i,\qquad y(Q_0)=s^2,
                              \quad 2Q_0=T_0.          \tag{A.82.6}
\]
Choose a function \(a_0'\) on \(E\) with divisor
\(T_0+O-2Q_0\). The divisor of \(a_0'\circ\phi\) is
\(D_K-2Q-2(Q+T)\), exactly the divisor of \(a_0a_2\).
After the finite constant extension needed to remove their ratio,
the opposite quotient is the actual connected normalization of
\(E_4\times_E D_0\), where
\[
 D_0:v^2=(z^2+q_0)^3-(z^2+q_0),\qquad
       q_{D_0}(z,v)=(z^2+q_0,v)+Q_0.                  \tag{A.82.7}
\]
Connectedness also follows from its already nontrivial character
and the computed branch divisor of size four.

The unramified double isogeny \(\phi\) is represented on its
target by the square class \(x-1\). This is literal:
\[
       x(\phi(P))-1=\left(\frac{y(P)}{2x(P)}\right)^2.
                                                               \tag{A.82.8}
\]
Its nonzero dual-kernel point is \((1,0)\), also the image of
the other source two-torsion points. In the fibre product replace
\(P\in E_4\) by \(P-Q\). Since \(\phi(Q)=Q_0\), this
removes the translation in (A.82.7). Therefore the unramified cover
of \(D_0\) is exactly obtained by adjoining
\(b\) with \(b^2=z^2+q_0-1\). Put \(c=v/b\). Its
three quadratic quotient fields over the \(z\)-line are
\[
 \begin{aligned}
 b^2&=z^2+q_0-1,\\
 c^2&=(z^2+q_0)(z^2+q_0+1),\\
 v^2&=(z^2+q_0-1)(z^2+q_0)(z^2+q_0+1).
 \end{aligned}                                        \tag{A.82.9}
\]
They have genera zero, one and two. Their actual \(C_2^2\)
projectors give \(J_{Y_{\{0,2\}}}\sim J_{D_0}\times F_*\),
where \(F_*\) is the middle quartic. The two involutions of
\(D_0\) give
\[
 J_{D_0}\sim E\times F_{128},\qquad
 F_{128}:y^2=x^3-4x^2+6x-4,\qquad j(F_{128})=128.
                                                               \tag{A.82.10}
\]
The old elliptic subspace from \(E_4\) equals that from \(E\),
because their pullbacks through the actual fibre-product square
agree and \(\phi\) is an isogeny. Consequently the opposite-pair
Prym is \(F_{128}\times F_*\), up to isogeny.
For the second opposite pair cyclic translation gives the same
geometric isogeny factors.

The elliptic quartic \(F_*\) has
\[
       j(F_*)=-11528-118792q_0.                        \tag{A.82.11}
\]
The following complete point counts certify that each of the three
new elliptic factors is non-CM. The final column is the squarefree
radicand of the ordinary Frobenius endomorphism field.
\[
 \begin{array}{c|r|r|r|r}
 \text{factor}&p&\#F(\mathbb F_p)&a_p&d\\ \hline
 F_{10976}&5&8&-2&-1\\
 F_{10976}&7&12&-4&-3\\
 F_{128}&5&8&-2&-1\\
 F_{128}&7&4&4&-3\\
 F_*&5&4&2&-1\\
 F_*&13&12&2&-3
 \end{array}                                         \tag{A.82.12}
\]
For \(F_*\), take \(q_0=2\) at five and \(q_0=5\) at
thirteen. Every displayed reduction is good and ordinary. The
retained program lists every quadratic-character value, including
the two points at infinity on the quartic. The ordinary elliptic
endomorphism theorem, \cite[V.3.1, printed p.145]{fghSilvermanAEC}, identifies the two fields.
They remain the same after every residue extension: the two
Frobenius eigenvalues have unequal valuations, so their ratio is
not a root of unity and no positive power becomes rational.
A hypothetical geometric CM field would specialize into both
quadratic fields after defining all endomorphisms over a finite
extension, an impossibility. This is the same persistent-field
argument used in A.80.

We have proved the actual geometric isogeny decomposition
\[
 J_C\sim E_4\times F^4\times(F')^4\times F_{10976}^4
                              \times(F_{128}\times F_*)^2.
                                                               \tag{A.82.13}
\]
Every factor except the displayed original elliptic factor is
non-CM. No assertion of mutual nonisogeny among those non-CM
factors is needed.

\paragraph{Exactly one main component in the actual special cover.}
Work at a place above two, after a finite extension making the
elliptic models good, the actual maps and marked points defined,
and the relevant curves semistable. A.78 gives potentially good
\(J_C\), special two-rank zero, and a compact-type proper stable
model. Fix one extension to \(k(C)\) of the vertical valuation of
the good elliptic model, and let \(D\le G\) be its decomposition
group. The intermediate isogeny \(E_4\to E\) has a two-primary
kernel and supersingular good special fibres. Its special map is
radicial, so it has one valuation over the base vertical valuation.
Consequently \(D\to G/N=C_4\) is onto.

The all-four Kummer character is the connected unramified double
cover of \(E_4\) whose class is the nonzero point \(T\) in
(A.82.3). Choose an origin after a finite extension. It is an actual
elliptic two-isogeny, again with good supersingular special fibres
and radicial special map. Thus its vertical valuation is unique,
and \(D\cap N\) surjects onto the total-parity quotient \(C_2\)
of \(N\). In particular it contains an odd-parity vector.
Conjugation by a lift in \(D\) of the cyclic generator preserves
\(D\cap N\) and acts on \(N\) by cyclic rotation. Identifying
\(N\) with \(\mathbb F_2[t]/((t+1)^4)\), every odd-parity
vector is a unit. Its rotations span \(N\). Hence
\[
                         D\cap N=N,\qquad D=G.         \tag{A.82.14}
\]
There is exactly one main component \(M\), meaning a component
whose map to the elliptic special fibre is nonconstant.

For completeness, this component assertion concerns the actual
model. The map extends through the relative Abel map and the
homomorphism of good Jacobians, as in
\cite[Theorem 1.1 and Proposition 3.2]{fghDupuyRabinoff}. Every nonconstant positive-genus component
would define an extension of the base vertical valuation; (A.82.14)
allows just one. Rational components cannot map nonconstantly to
a smooth elliptic curve. Thus all other components are contracted.
A.78's main-component comparison identifies the normalization of
\(M\), up to a radicial map, with the connected separable
\(2\)-group cover of the elliptic special fibre obtained from
its decomposition and inertia groups away from \(O\). It is
branched at most at \(O\). Both source and base have two-rank
zero. Deuring--Shafarevich, \cite[Corollary 1.8]{fghCrew1984}, gives
\(-1=-\#\{\text{points above }O\}\), including the unramified
case, so the normalization of \(M\) has just one point above
\(O\).

The same conclusions hold for every actual intermediate \(Z\).
Indeed the decomposition group \(D=G\) is transitive on every
coset set \(G/H\), so its vertical valuation has one extension
to \(k(Z)\). The finite residue-field map from the normalization
of \(M\) to the intermediate main normalization is surjective;
the latter also has only one point above \(O\). The norm
\(J_C\to J_Z\) is surjective, so \(J_Z\) has good reduction
and two-rank zero and its semistable graph is a tree.

\paragraph{The precise positive distance supplied by the marks.}
The base quotient \(Z=E\) is already covered by A.40.
Suppose \(H\ne G\). Inertia in \(C/E\) has order two.
Its orbits on a degree \(d=[G:H]>2\) coset set give at least
\(d/2\ge2\) cusps. If \(d=2\), the inertia commutator is
killed in the abelian quotient, so the cover is unramified and
has two cusps. Thus every selected cusp has a distinct companion.
For genus-one \(Z\), A.40 already gives the conclusion; otherwise
use its proper stable model and then separate all marked sections.

Every cusp maps to \(O\). All contracted components in a connected
component of the graph minus the main vertex have the same image
in the elliptic curve. A component carrying a cusp therefore has
image \(O\), and is attached to the unique point of the main
normalization above \(O\). If marks also specialize at that point
on the main component, separating the distinct sections inserts
a rational tree there. Since a semistable node has two branches,
there is a single attached subtree carrying all the marks over
that point. In the marked model the path between any two distinct
cusps therefore avoids the main positive-genus vertex. In
particular, after choosing the companion \(Q_0\),
\[
 r_M=\operatorname{dist}(M,\text{the }P\text{--}Q_0
                              \text{ path})>0.         \tag{A.82.15}
\]
The path may pass through other positive-genus components. No
rational-path assertion for the whole path is being made.

\paragraph{The partial-distance kernel for a main block of any genus.}
At an odd coefficient prime, let \(V_M\) be the local homology
block of the main component, and let \(V_{\rm tail}\) be the
sum of all other handle blocks. These are local subspaces only.
Root at \(P\), eliminate \(t_{Q_0}\), and retain the remaining
cusp letters \(T\), of weight two; handles have weight one.
The genuine first weighted edge derivation of A.41 has
\(Dv=[W_i,v]\), \(Dt_j=[W_j,t_j]\), with conjugators linear
in the component bivectors and cusp letters. Its main self
coefficient is
\[
            [\Omega_M]W_M
              =\sum_e m_e M_e(1-Q_{0,e})=r_M>0.       \tag{A.82.16}
\]
The arbitrary packet term \(\lambda\operatorname{ad}(t_P)\)
does not change this coefficient. The actual local Frobenius--tame
commutator supplies a nonzero common multiple of this derivation
inside the closed derived Galois group. The based comparison and
absence of a weight-one inner discrepancy are precisely A.41.4--6.

Set \(B=\mathbb Q_\ell\langle T,V_{\rm tail},\Omega_M\rangle\).
This is a free algebra on the displayed primitive generators.
Use A.41.8 only for the main handle algebra: its leading
symplectic word has coefficient one and no self-overlap, for
any genus of \(M\). All other handles remain coefficient letters.
The resulting normal basis has a first nonboundary main core
preceded by a \(B\)-word. The derivation preserves \(B\).
On a fixed first-core tail, the top first-slot operator is
\(D_B+R_{W_M}\), with first-slot size measured by weight.
Later-core terms cannot contribute to that top slot.

Count occurrences of \(\Omega_M\) in the slot. Its increasing
part sends any other generator \(z\) to
\(c_z[\Omega_M,z]\), and contributes \(r_M R_{\Omega_M}\).
For a fixed word \(z_1\cdots z_b\) in the other generators,
record the powers of \(\Omega_M\) in its \(b+1\) slots by
commuting variables \(Z_0,\ldots,Z_b\). The top operator is
multiplication by
\[
 c_{z_1}Z_0+\sum_{j=1}^{b-1}(c_{z_{j+1}}-c_{z_j})Z_j
                              +(r_M-c_{z_b})Z_b.       \tag{A.82.17}
\]
Its coefficients sum to \(r_M\ne0\); for an empty other word
it is \(r_MZ_0\). Hence it is injective. This proves, in every
weight and with no finite-degree extrapolation,
\[
             \ker D\subseteq
                \mathbb L(T,V_{\rm tail},\Omega_M).   \tag{A.82.18}
\]
The proof is the actual A.80 partial-distance argument with all
nonmain handles retained. It uses no positivity for other blocks.

\paragraph{The global comparison uses only the actual elliptic map.}
Let \(V_E\) be the image of the actual pullback from \(E\)
in \(V_\ell J_Z\), and put \(W=\ker(f_{Z,*})\).
Polarization adjunction gives the global orthogonal splitting
\(V_\ell J_Z=V_E\oplus W\), since
\(f_{Z,*}f_Z^*=[\deg f_Z]\) is invertible over \(\mathbb Q_\ell\).
Specialization of the actual norm map is zero on every contracted
component and nonzero on \(M\). Thus \(V_E\subset V_M\) and
\(V_{\rm tail}\subset W\). Inside the local main block choose
its orthogonal splitting
\[
 V_M=V_E\oplus W_M,\quad W_M\subset W,\qquad
                  \Omega_M=c\Omega_E+\Omega_{W_M},
                              \quad c\ne0.            \tag{A.82.19}
\]
The scalar is the nonzero polarization scale on the actual
elliptic pullback; the main bivector uses the restricted ambient
symplectic form. No global stability of \(V_M\) or of \(W_M\)
is asserted. Since \(\Omega_{W_M}\in\Lambda^2W\), (A.82.18)
implies the larger, globally stable bound
\[
                    \ker D\subseteq\mathbb L(T,W,\Omega_E).
                                                               \tag{A.82.20}
\]

The actual norm \(J_C\to J_Z\) makes \(J_Z\) an isogeny
factor of (A.82.13). It contains the original elliptic factor,
and there is only one copy of that CM isogeny class in (A.82.13).
Consequently \(W\) is a sum of Tate modules of non-CM elliptic
curves, with multiplicities. This is not a claim that the local
main block is a global isogeny factor.

\paragraph{Product monodromy and the actual fixed subgroup.}
After a finite extension defining all geometric isogenies, group
the non-CM elliptic factors by geometric isogeny class. Serre's
Theorem 3, printed p.299, \cite{A_Serre1972}, gives full
\(\mathfrak{gl}_2\) projection on each class.
Faltings, Section 5, Satz 3, \cite{DE_Faltings1983}, gives reductivity of the
connected monodromy group. Its derived Lie algebra is a semisimple
subdirect product of the \(\mathfrak{sl}_2\)'s. Lie-Goursat says
that a proper such product has a diagonal correlation between two
factors. Every automorphism of split \(\mathfrak{sl}_2\) is inner.
The resulting intertwiner also commutes with the central Lie
algebra: equal determinants give equal scalar actions. It is then
fixed by the connected group, hence by an open Galois subgroup.
Faltings' Section 5, Korollar 1, printed p.361, \cite{DE_Faltings1983},
makes it an actual
nonzero homomorphism between the two elliptic curves, contradicting
geometric nonisogeny. The connected derived group is therefore the
full product of \(\mathrm{SL}_2\)'s, with identical standard
representations on repeated isogeny factors.

In its invariant tensors, each class has even total degree, since
its central \(-I\) acts trivially. On such an invariant of degree
\(2b\), any actual Galois matrix acts by the common determinant
\(\chi_\ell^b\): over an algebraic closure factor each matrix
as a scalar times an \(\mathrm{SL}_2\) matrix. The ambiguity in
the square root disappears in even degree. The same conclusion
holds for invariant subquotients by reductivity. Both \(T\) and
\(\Omega_E\) have weight two and character \(\chi_\ell\).
Thus every invariant subspace of (A.82.20) has no odd weight and
has scalar character \(\chi_\ell^b\) in weight \(2b\).

Now take the actual subgroup fixed by the closed derived group
of a sufficiently small number-field Galois group. Its finite
weighted unipotent closures are stable and have abelian arithmetic
action. Their leading classes are killed by the genuine local
commutator derivation, hence satisfy (A.82.20), and are invariant
under the connected derived monodromy group. The preceding
calculation gives exactly the Tate-graded hypothesis of A.38.
That actual-subgroup criterion proves (A.82.1) at every odd prime,
without lifting arbitrary rational Lie elements. The selection of
a further finite field places its closed derived group inside any
prescribed open derived subgroup, exactly as in A.38/A.80.
At two, A.22 applies to every intermediate of this actual two-group
closure for the original packets. A.59 gives all stated normalizer
assertions. Arithmetic identifications use actual maps and are
inner-invariant; no profinite geometric path is declared rational
at a finite field extension.

The complete exact program \path{notes/fgh_wreath64_jacobian_probe.py}
retains the divisor vectors, all fifteen character dimensions,
the quotient identities and every finite-field character value.
The actual model, kernel and monodromy proofs above are separate
from those finite checks. This candidate remains unregistered
pending complete independent and root audits.
\end{proof}
